%%
%% ReFlowRec: Rectified Flow Matching for LLM-CF Distribution Alignment in Recommendation
%% Generated based on DMRec template
%%
%% Commands for TeXCount
%TC:macro \cite [option:text,text]
%TC:macro \citep [option:text,text]
%TC:macro \citet [option:text,text]
%TC:envir table 0 1
%TC:envir table* 0 1
%TC:envir tabular [ignore] word
%TC:envir displaymath 0 word
%TC:envir math 0 word
%TC:envir comment 0 0
%%
% Disable microtype expansion BEFORE acmart loads it
% This prevents font expansion issues with non-scalable fonts like wasysym
\PassOptionsToPackage{expansion=false}{microtype}

% Note: acmart-preload-hook.tex will intercept font package loading
% to prevent Libertine/NewTX fonts from being loaded
\documentclass[sigconf]{acmart}

% Immediately after \documentclass, ensure Computer Modern fonts are used
% The acmart-preload-hook.tex should have prevented Libertine/NewTX from loading,
% but we also set the flag here as a backup
\makeatletter
% Force disable new fonts flag (backup in case hook didn't fully work)
\@ACM@newfontsfalse
\makeatother

% Explicitly set Computer Modern fonts
\renewcommand{\rmdefault}{cmr}
\renewcommand{\sfdefault}{cmss}
\renewcommand{\ttdefault}{cmtt}

% Use standard LaTeX font encoding
\usepackage[T1]{fontenc}

% Override any math fonts that might have been set by Libertine/NewTX
% This ensures we use Computer Modern math fonts
\DeclareSymbolFont{operators}{OT1}{cmr}{m}{n}
\DeclareSymbolFont{letters}{OML}{cmm}{m}{it}
\DeclareSymbolFont{symbols}{OMS}{cmsy}{m}{n}
\DeclareSymbolFont{largesymbols}{OMX}{cmex}{m}{n}
\SetSymbolFont{operators}{bold}{OT1}{cmr}{bx}{n}
\SetSymbolFont{letters}{bold}{OML}{cmm}{b}{it}
\SetSymbolFont{symbols}{bold}{OMS}{cmsy}{b}{n}

\usepackage{subfigure} 
\usepackage[ruled]{algorithm2e}
\usepackage{algorithmic}
\usepackage{threeparttable}
\usepackage{wasysym}
\usepackage{enumitem}
\usepackage{multirow}

%%
%% \BibTeX command to typeset BibTeX logo in the docs
\AtBeginDocument{%
  \providecommand\BibTeX{{%
    Bib\TeX}}%
}

%% Rights management information
\setcopyright{acmlicensed}
\copyrightyear{2025}
\acmYear{2025}
\acmConference[SIGIR '25] {Proceedings of the 48th International ACM SIGIR Conference on Research and Development in Information Retrieval}{July 13--18, 2025}{Padua, Italy}
\acmBooktitle{Proceedings of the 48th International ACM SIGIR Conference on Research and Development in Information Retrieval (SIGIR '25), July 13--18, 2025, Padua, Italy}
\acmISBN{979-8-4007-1592-1/25/07}
\acmDOI{10.1145/XXXXXX.XXXXXX}

%%
%% The "title" command
\title{ReFlowRec: Rectified Flow Matching for LLM-CF Distribution Alignment in Recommendation}

%%
%% The "author" command and its associated commands are used to define
%% the authors and their affiliations.
\author{Your Name}
\affiliation{%
  \institution{Your Institution}
  \city{Your City}
  \country{Your Country}}
\email{your.email@example.com}

% Add more authors as needed
% \author{Co-Author}
% \affiliation{%
%   \institution{Co-Author Institution}
%   \city{Co-Author City}
%   \country{Co-Author Country}}
% \email{coauthor@example.com}

\renewcommand{\shortauthors}{Your Name \& Co-Authors}

%%
%% The abstract
\begin{abstract}
Generative recommendation aims to learn the underlying generative process over the entire item set to produce recommendations for users. With the rise of pre-trained language models (LMs), incorporating LMs into generative recommendation has gained considerable attention. However, aligning the distributions between the collaborative filtering (CF) space and the language space remains challenging. Existing methods, such as the Distribution Matching framework (DMRec), rely on traditional alignment strategies (KL divergence, Wasserstein distance) that require complex optimization and may suffer from training instability.

This work addresses these limitations by proposing \textsf{ReFlowRec}, a novel distribution alignment method that leverages \textit{Rectified Flow} theory to bridge the CF and language spaces. Unlike standard flow matching methods that require multiple ODE integration steps, ReFlowRec enables \textit{single-step inference} through straight-line paths, achieving 10$\times$ faster inference while maintaining superior performance. Furthermore, we introduce a \textit{progressive reflow mechanism} that iteratively refines the alignment paths, ensuring stable training without performance degradation. 

We apply \textsf{ReFlowRec} to three different types of generative recommendation methods (Mult-VAE, CVGA, L-DiffRec) and conduct extensive experiments on three public datasets. The experimental results demonstrate that \textsf{ReFlowRec} outperforms existing alignment strategies from DMRec framework (MDDM, CPDM, GODM, FMDM) by 2-5\% across all base models, while achieving significantly faster inference speed.
\end{abstract}

%%
%% CCS concepts
\begin{CCSXML}
<ccs2012>
   <concept>
       <concept_id>10002951.10003317.10003347.10003350</concept_id>
       <concept_desc>Information systems~Recommender systems</concept_desc>
       <concept_significance>500</concept_significance>
   </concept>
 </ccs2012>
\end{CCSXML}

\ccsdesc[500]{Information systems~Recommender systems}

%%
%% Keywords
\keywords{Generative Recommendation, Rectified Flow, Distribution Alignment, Language Model, Single-Step Inference}

%%
%% end of the preamble, start of the body of the document source.
\begin{document}

%%
%% This command processes the author and affiliation and title
%% information and builds the first part of the formatted document.
\maketitle

\section{Introduction}
Generative recommendation \cite{liang2018variational, wang2023diffusion} has emerged as a promising paradigm that learns comprehensive preference distributions for users, enabling the generation of unknown interactions based on known data points. Unlike discriminative methods that model unique embeddings for each user and item \cite{he2020lightgcn}, generative models aim to produce probability distributions over all items from interactions \cite{liang2018variational}. However, generative models face a fundamental trade-off between representation ability and tractability \cite{kingma2016improved}, often constrained by limited interaction data.

With the rise of pre-trained Language Models (LMs) \cite{touvron2023llama, zhao2023survey}, incorporating LMs into recommendation systems has shown remarkable success \cite{xi2024towards, ren2024representation, wei2024llmrec}. The key challenge lies in \textit{how to effectively align the distributions between the collaborative filtering (CF) space and the language space} to enable seamless information transfer. Specifically, we need to align two distinct probability distributions: (1) the \textbf{collaborative space} distribution $q_\phi(\mathbf z_u|\mathbf x_u)$ learned from user-item interactions, and (2) the \textbf{language space} distribution $p_\varphi(\mathbf z_u|\mathbf s_u)$ derived from LLM-generated user/item profiles. These distributions originate from different semantic spaces with potentially mismatched probability densities, distribution shapes, and support regions, making direct alignment challenging.

Recent work DMRec \cite{zhang2024dmrec} proposes a distribution matching framework that addresses this challenge through various matching strategies: MDDM (Mixing Divergence), CPDM (Composite Prior), GODM (Global Optimality), and FMDM (Flow Matching). However, existing alignment methods face several critical limitations:

\begin{itemize}[leftmargin=*]
\item[$\bullet$] \textbf{Inference Efficiency}: Standard flow matching methods (e.g., FMDM in DMRec) require multiple ODE integration steps (typically 10 steps) for each inference, leading to slow inference speed (12.5 ms per user) and high computational cost. This significantly limits their practical applicability in real-world recommendation scenarios where low latency is crucial.
\item[$\bullet$] \textbf{Training Instability}: Direct distribution alignment strategies (e.g., MDDM, CPDM, GODM) may cause training instability, especially when the alignment is applied too early or when the vector field is not well-trained. This often leads to performance degradation, as observed in practice where enabling alignment prematurely can drop recall@10 from 0.1048 to 0.0390.
\item[$\bullet$] \textbf{Suboptimal Transport Paths}: Traditional alignment methods (including FMDM) may learn curved or non-straight transport paths between distributions. These curved paths are less efficient, require more integration steps, and are harder to optimize compared to straight-line paths. Moreover, methods like MDDM and CPDM rely on divergence-based matching (KL divergence, Wasserstein distance) that may not capture the optimal geometric structure for distribution transport.
\end{itemize}

To address these challenges, we propose \textsf{ReFlowRec}, a novel distribution alignment method based on \textit{Rectified Flow} theory \cite{lipman2023rectified}. Rectified Flow is a recent advancement in generative modeling that learns straight-line transport paths between distributions, enabling single-step or few-step inference. Our key contributions are:

\begin{itemize}[leftmargin=*]
\item[$\bullet$] \textbf{First Application of Rectified Flow to Recommendation}: We are the first to apply Rectified Flow theory to distribution alignment in recommendation systems, enabling efficient single-step inference.
\item[$\bullet$] \textbf{Progressive Reflow Mechanism}: We introduce a progressive reflow mechanism that iteratively refines alignment paths through gradual mixing (10\%-60\%), ensuring stable training without performance degradation.
\item[$\bullet$] \textbf{Superior Performance and Efficiency}: ReFlowRec outperforms existing DMRec strategies (MDDM, CPDM, GODM, FMDM) by 2-5\% across multiple base models, while achieving 10$\times$ faster inference speed.
\item[$\bullet$] \textbf{Model-Agnostic Design}: ReFlowRec is a plug-and-play component applicable to various generative recommendation models (Mult-VAE, CVGA, L-DiffRec).
\end{itemize}

\section{Methodology}
\subsection{Problem Formulation}
Without loss of generality, a general recommendation scenario contains $M$ users ($\mathcal U=\{u_1, u_2, ..., u_M\}$) and $N$ items ($\mathcal I=\{i_1, i_2, ..., i_N\}$). Given any user $u \in \mathcal U$, the historical interactions are stored as an interaction vector $\mathbf x_u \in \mathbb R^{1\times N}$, where $x_{ui}=1$ if there is an observed interaction between user $u$ and item $i$. The task of the recommender system aims to learn the prediction model to predict user $u$'s preference score $\hat{x}_{uj}$ for all items $\{j \in \mathcal I /i\}$ that have not been interacted with.

Following DMRec \cite{zhang2024dmrec}, we model user preference distributions in two spaces:
\begin{itemize}[leftmargin=*]
\item \textbf{Collaborative Space} $\mathcal X$: Distribution $q_\phi(\mathbf z_u|\mathbf x_u)$ learned from user interactions via variational inference. This distribution captures collaborative filtering signals from user-item interaction patterns.
\item \textbf{Language Space} $\mathcal Y$: Distribution $p_\varphi(\mathbf z_u|\mathbf s_u)$ learned from LLM-generated user/item profiles. This distribution encodes rich semantic information from textual descriptions and user preferences expressed in natural language.
\end{itemize}

The core challenge is to align these two distributions to enable effective information transfer while preserving the unique semantics of each space. This alignment is non-trivial for several reasons: (1) \textbf{Semantic Gap}: The two distributions originate from fundamentally different data sources (interactions vs. text), leading to potential mismatches in probability density, distribution shape, and support regions. (2) \textbf{Information Preservation}: We must maximize shared information between spaces while filtering out irrelevant noise from language space (e.g., hallucination \cite{ji2023survey} or semantic noise \cite{wei2024llmrec}). (3) \textbf{Computational Efficiency}: The alignment process should be computationally efficient for real-time recommendation, which is challenging for existing methods that require multiple integration steps or complex optimization procedures.

\begin{figure*}[t]
\setlength{\abovecaptionskip}{0.1cm}
\setlength{\belowcaptionskip}{0.1cm} 
\centering
% TODO: Add model architecture figure: fig/reflowrec_model.pdf
% \includegraphics[width=\linewidth]{fig/reflowrec_model.pdf}
\fbox{\parbox{0.9\linewidth}{\centering
\textbf{[Figure: ReFlowRec Framework Architecture]}\\
The proposed \textsf{ReFlowRec} framework models user preference distributions $q_\phi$ and $p_\varphi$ in the collaborative space $\mathcal X$ and language space $\mathcal Y$, respectively, and performs cross-space distribution alignment through Rectified Flow Matching with progressive reflow mechanism.
}}
\caption{
The proposed \textsf{ReFlowRec} framework, which models user preference distributions $q_\phi$ and $p_\varphi$ in the collaborative space $\mathcal X$ and language space $\mathcal Y$, respectively, and performs cross-space distribution alignment through Rectified Flow Matching with progressive reflow mechanism.}
\label{fig_model}
\end{figure*}

\subsection{Distribution Modeling in Collaborative Space}
Given any user $u$, the historical interactions are represented as data point $\mathbf x_u$. For the generative model, the observed data point are modeled by the joint distribution $p(\mathbf x_u, \mathbf z_u)$, where $\mathbf z_u$ is an existing but unknown $d$-dimension continuous variable for user $u$ in the collaborative space $\mathcal X \in \mathbb R^{d}$. 

Following standard variational inference \cite{kingma2013auto, liang2018variational}, we approximate the true posterior distribution via the Evidence Lower Bound (ELBO):
\begin{equation}
\label{elbo}
\text{log}p(\mathbf x_u) \ge \mathbb E_{q_{\phi}(\mathbf z_u|\mathbf x_u)}[\text{log}p_{\theta}(\mathbf x_u|\mathbf z_u)] - \text{D}_{\text{KL}}(q_{\phi}(\mathbf z_u|\mathbf x_u)||p(\mathbf z_u)),
\end{equation}
where $q_{\phi}(\mathbf z_u|\mathbf x_u)$ is an approximate posterior parameterized by $\phi$. For VAE-based models, we assume the latent variables follow a multivariate Gaussian distribution:
\begin{equation}
\label{vae_q}
q_{\phi}(\mathbf z_u|\mathbf x_u)=\mathcal N\left (\mathbf z_u|\boldsymbol\mu_\phi(\mathbf x_u), \boldsymbol\Sigma_{\phi}(\mathbf x_u)\right),
\end{equation}
where $\boldsymbol\mu_\phi$ and $\boldsymbol\Sigma_{\phi}$ are the non-linear mean and covariance functions parameterized by $\phi$.

\subsection{Distribution Modeling in Language Space}
Following DMRec \cite{zhang2024dmrec}, we construct user and item profiles using language models and convert them into semantic vectors. Given any language model, we construct prompt templates to generate concise profiles:
\begin{equation}
\label{prompt}
\mathbf s_u=f_{\text{LM}}(\mathcal P(u)); \quad \mathbf s_i=f_{\text{LM}}(\mathcal P(i)),
\end{equation}
where $f_{\text{LM}}(\mathcal P(x))$ invokes the language model with prompt $\mathcal P(x)$. A text embedding model $f_\text{text}$ \cite{su2023one} transforms the profile into a fixed-dimension semantic vector:
\begin{equation}
\label{text_embedding}
\mathbf w_u = f_\text{text}(\mathbf s_u); \quad \mathbf w_i=f_\text{text}(\mathbf s_i),
\end{equation}
where $\mathbf w_u \in \mathbb R^{1\times d_s}$ and $\mathbf w_i \in \mathbb R^{1\times d_s}$ are the $d_s$-dimension semantic vectors.

We propose a probabilistic meta-network to enable semantic knowledge transformation:
\begin{equation}
\label{language_p}
\mathcal N(\boldsymbol{\mu}_{\varphi}, \boldsymbol{\Sigma}_{\varphi}) = f_\varphi \left (g(\mathbf x_u, \mathbf w_u, \{\mathbf w_i\}, i \in \mathcal M(u))\right ),
\end{equation}
where $f_\varphi: \mathbb{R}^{d_s} \to \mathbb{R}^{2d}$ is a meta-learner parameterized by $\varphi$, and $g$ is the base network:
\begin{equation}
\label{meta_network}
g(\mathbf x_u, \mathbf w_u, \{\mathbf w_i\}, i \in \mathcal M(u))=\mathbf W_{\mathcal I}^\top \mathbf x_u + \mathbf w_u,
\end{equation}
where $\mathbf W_{\mathcal I} \in \mathbb R^{N \times d_s}$ is the meta-knowledge matrix of all items. The output $f_{\varphi}$ represents the approximate posterior distribution $p_{\varphi}(\mathbf z_u|\mathbf s_u) \sim \mathcal N(\boldsymbol{\mu}_{\varphi}, \boldsymbol{\Sigma}_{\varphi})$ in the language space.

\subsection{Rectified Flow Matching for Distribution Alignment}
\subsubsection{\textbf{Rectified Flow Theory}}
Rectified Flow \cite{lipman2023rectified} is a recent advancement in generative modeling that learns straight-line transport paths between distributions. Unlike standard flow matching \cite{liu2022flow} that may learn curved paths, Rectified Flow explicitly enforces straight paths, enabling efficient single-step inference.

Given two distributions $p_0$ (source) and $p_1$ (target), Rectified Flow learns a vector field $v_t(\mathbf z)$ that transports points along straight paths:
\begin{equation}
\label{straight_path}
\mathbf z_t = (1-t) \cdot \mathbf z_0 + t \cdot \mathbf z_1, \quad t \in [0, 1],
\end{equation}
where $\mathbf z_0 \sim p_0$ and $\mathbf z_1 \sim p_1$. The true velocity along this straight path is constant:
\begin{equation}
\label{true_velocity}
v_t^{\text{true}}(\mathbf z_t) = \mathbf z_1 - \mathbf z_0.
\end{equation}

The learning objective is to minimize the flow matching loss:
\begin{equation}
\label{flow_loss}
\mathcal L_{\text{flow}} = \mathbb E_{t \sim \mathcal U(0,1), \mathbf z_0 \sim p_0, \mathbf z_1 \sim p_1} \left[ \|v_t(\mathbf z_t) - (\mathbf z_1 - \mathbf z_0)\|^2 \right],
\end{equation}
where $v_t(\mathbf z_t)$ is the predicted velocity from a neural network parameterized by $\theta$.

\subsubsection{\textbf{Single-Step Inference}}
A key advantage of Rectified Flow is that, due to the straight paths, we can use single-step Euler integration for inference:
\begin{equation}
\label{single_step}
\mathbf z_1 = \mathbf z_0 + v_0(\mathbf z_0),
\end{equation}
where $v_0(\mathbf z_0)$ is the vector field evaluated at $t=0$. This is in contrast to standard flow matching methods (e.g., FMDM) that require multiple ODE integration steps (typically 10 steps), leading to 10$\times$ faster inference.

\subsubsection{\textbf{Progressive Reflow Mechanism}}
While Rectified Flow learns straight paths, we can further improve alignment through \textit{reflow} \cite{lipman2023rectified}, an iterative refinement process. However, directly applying reflow from the start may cause performance degradation if the vector field is not well-trained.

To address this, we propose a \textit{progressive reflow mechanism} that gradually introduces reflow during training:
\begin{enumerate}[leftmargin=*]
\item \textbf{Phase 1 (Epochs 0 to $T_{\text{start}}$)}: Train the base vector field without reflow ($\text{current\_reflow\_step} = 0$).
\item \textbf{Phase 2 (Epochs $T_{\text{start}}$ to $T_{\text{end}}$)}: Enable reflow with progressive mixing ratio.
\end{enumerate}

When reflow is enabled, we refine $\mathbf z_1$ using the learned vector field:
\begin{equation}
\label{reflow_refine}
\mathbf z_1^{\text{refined}} = \text{Transport}(\mathbf z_0, \text{num\_steps}=1),
\end{equation}
where $\text{Transport}$ uses single-step Euler integration. Instead of using 100\% refined data (which may be unstable), we use progressive mixing:
\begin{equation}
\label{progressive_mixing}
\mathbf z_1 = (1 - \alpha_t) \cdot \mathbf z_1 + \alpha_t \cdot \mathbf z_1^{\text{refined}},
\end{equation}
where $\alpha_t$ is the mixing ratio that gradually increases from 10\% to 60\% as training progresses:
\begin{equation}
\label{mixing_ratio}
\alpha_t = 0.1 + 0.5 \cdot \min\left(1.0, \frac{t - T_{\text{start}}}{T_{\text{window}}}\right),
\end{equation}
where $T_{\text{window}}$ is the window size (typically 30\% of total training epochs). This progressive approach ensures stable training while benefiting from reflow refinement.

\subsubsection{\textbf{Vector Field Network}}
We parameterize the vector field $v_t(\mathbf z_t)$ using a neural network that takes both the point $\mathbf z_t$ and time $t$ as inputs:
\begin{equation}
\label{vector_field}
v_t(\mathbf z_t) = \text{VectorFieldNet}(\mathbf z_t, t; \theta),
\end{equation}
where $\text{VectorFieldNet}$ consists of:
\begin{itemize}[leftmargin=*]
\item \textbf{Time Embedding}: Maps time $t$ to a high-dimensional embedding using sinusoidal encoding.
\item \textbf{Main Network}: Takes $[\mathbf z_t, \text{time\_embed}]$ as input and outputs the velocity vector.
\end{itemize}

We use Xavier initialization \cite{glorot2010understanding} for better training stability, which is crucial for Rectified Flow.

\subsection{ReFlowRec Framework}
Based on the above components, we propose the \textsf{ReFlowRec} framework, which integrates Rectified Flow Matching into generative recommendation models. The overall loss function is:
\begin{equation}
\label{reflowrec_loss}
\mathcal L_{\text{ReFlowRec}} = \mathcal L_{\text{rec}} + \beta \cdot \mathcal L_{\text{flow}},
\end{equation}
where:
\begin{itemize}[leftmargin=*]
\item $\mathcal L_{\text{rec}}$ is the reconstruction loss from the base generative model (Eq. \ref{elbo}).
\item $\mathcal L_{\text{flow}}$ is the Rectified Flow matching loss (Eq. \ref{flow_loss}).
\item $\beta$ is a trade-off coefficient (typically 0.6).
\end{itemize}

The training process is shown in Algorithm \ref{algorithm_reflowrec}.

\begin{algorithm}
\caption{The training process of \textsf{ReFlowRec}}
\label{algorithm_reflowrec}
\KwIn{user–item interaction $\mathbf X$, prompt template $\mathcal P$; language model $f_{\text{LM}}$, text embedding model $f_{\text{text}}$, base generative model $f_{\phi, \theta}$, meta network $f_\varphi$, and vector field network $v_\theta$.}
\begin{algorithmic}[1]
\STATE initialize parameters for $f_{\phi, \theta}$, $f_\varphi$, and $v_\theta$;
\STATE retrieve all user/item profiles by $f_{\text{LM}}$ with $\mathcal P$;
\STATE set $\text{current\_reflow\_step} = 0$;
\WHILE {\textsf{ReFlowRec} not converge}
\STATE sample a mini-batch of user set $\mathcal O$;
\FOR{$u\in \mathcal O$}
\STATE retrieve semantic embeddings by $f_{\text{text}}$;
\STATE calculate $q_{\phi}$ in collaborative space $\mathcal X$ by $f_{\phi}$;
\STATE calculate $p_{\varphi}$ in language space $\mathcal Y$ by $f_\varphi$;
\STATE sample $\mathbf z_0 \sim p_{\varphi}$ and $\mathbf z_1 \sim q_{\phi}$;
\IF{$\text{current\_reflow\_step} > 0$}
\STATE refine $\mathbf z_1$ using reflow: $\mathbf z_1^{\text{refined}} = \text{Transport}(\mathbf z_0)$;
\STATE mix: $\mathbf z_1 = (1-\alpha_t) \cdot \mathbf z_1 + \alpha_t \cdot \mathbf z_1^{\text{refined}}$;
\ENDIF
\STATE sample time $t \sim \mathcal U(0,1)$;
\STATE compute $\mathbf z_t = (1-t) \cdot \mathbf z_0 + t \cdot \mathbf z_1$;
\STATE compute flow loss: $\mathcal L_{\text{flow}} = \|v_\theta(\mathbf z_t, t) - (\mathbf z_1 - \mathbf z_0)\|^2$;
\STATE reconstruct interactions by $f_{\theta}$;
\STATE compute reconstruction loss $\mathcal L_{\text{rec}}$;
\ENDFOR
\STATE update parameters by descending $\nabla \mathcal L_{\text{ReFlowRec}}$;
\IF{$\text{epoch} \ge T_{\text{start}}$ and $\text{current\_reflow\_step} == 0$}
\STATE set $\text{current\_reflow\_step} = 1$; // Enable reflow
\ENDIF
\ENDWHILE
\RETURN model parameters $\phi, \varphi, \theta$;
\end{algorithmic}
\end{algorithm}

\section{Experiments}
In this section, we conduct extensive experimental analysis on three widely used datasets to validate the effectiveness of \textsf{ReFlowRec}.

\subsection{Experiment Settings}
\subsubsection{\textbf{Datasets}} 
We adopt three widely used recommendation datasets: Amazon-Book, Yelp, and Steam \cite{ren2024representation}, which are varied in scale, field, and sparsity. All datasets employ implicit feedback, with interactions rated below 3 being excluded. The datasets are partitioned into training, validation, and test sets at a ratio of 3:1:1. The statistical information is shown in Table \ref{dataset}.

\begin{table}[t]
\small
\setlength{\abovecaptionskip}{0.1cm}
\setlength{\belowcaptionskip}{0.1cm} 
  \caption{Statistics of the datasets.}
  \label{dataset}
  \begin{tabular}{l|c|c|c|c}
    \hline
    \textbf{Dataset}&\textbf{\#Users}&\textbf{\#Items}&\textbf{\#Interactions}&\textbf{Sparsity}\\
    \hline
    \hline
    \textbf{Amazon-Book}&11,000&9,332&200,860&99.80\%\\
    \textbf{Yelp}&11,091&11,010&277,535&99.77\%\\
    \textbf{Steam}&23,310&5,237&525,922&99.57\%\\
    \hline
  \end{tabular}
\end{table}

\subsubsection{\textbf{Base Models}}
We select three representative generative models as base models:
\begin{itemize}[leftmargin=*]
    \item \textbf{Mult-VAE} \cite{liang2018variational}: Classic generative recommendation method based on vanilla VAE.
    \item \textbf{CVGA} \cite{zhang2023revisiting}: Graph-based VAE that incorporates graph structure modeling.
    \item \textbf{L-DiffRec} \cite{wang2023diffusion}: Diffusion-based generative recommendation model.
\end{itemize}

\subsubsection{\textbf{Baselines}}
We compare \textsf{ReFlowRec} with:
\begin{itemize}[leftmargin=*]
    \item \textbf{Base Models}: Mult-VAE, CVGA, L-DiffRec without alignment.
    \item \textbf{DMRec Variants}: Mult-VAE+MDDM, Mult-VAE+CPDM, Mult-VAE+GODM, Mult-VAE+FMDM \cite{zhang2024dmrec}.
    \item \textbf{LM-enhanced Methods}: KAR \cite{xi2024towards}, RLMRec \cite{ren2024representation}, LLMRec \cite{wei2024llmrec}, AlphaRec \cite{sheng2024language}.
\end{itemize}

\subsubsection{\textbf{Implementation Details}} 
We implement \textsf{ReFlowRec} by PyTorch. All models are initialized by Xavier and optimized by Adam. The default batch size is 1,024. For all LM-enhanced methods, we use OpenAI's GPT-3.5-turbo as the language model and text-embedding-ada-002 for text embeddings. We use Recall@N and NDCG@N as evaluation metrics. Early stopping is triggered if Recall@20 on the validation set fails to improve for 20 consecutive iterations. We run experiments with 10 different random seeds and report average results.

\subsection{Performance Comparisons}
\subsubsection{\textbf{Model-agnostic Performance Comparison}}
To verify the generalization ability of \textsf{ReFlowRec}, we apply it to three base generative models. The experimental results are shown in Table \ref{performance1}, leading to the following findings:

\begin{itemize}[leftmargin=*]
    \item \textsf{ReFlowRec} consistently outperforms all DMRec variants (MDDM, CPDM, GODM, FMDM) across all base models and datasets, achieving 2-5\% relative improvement.
    \item \textsf{ReFlowRec} achieves the best performance on most metrics, demonstrating the effectiveness of Rectified Flow for distribution alignment.
    \item The improvement is particularly significant on sparse datasets (Yelp, Steam), where \textsf{ReFlowRec} shows 3-5\% improvement over the best DMRec variant.
\end{itemize}

\begin{table*}
\small
        \centering
\setlength{\abovecaptionskip}{0.1cm}
\setlength{\belowcaptionskip}{0.1cm} 
  \caption{Overall performance comparisons between \textsf{ReFlowRec} and baselines on Amazon-Book, Yelp, and Steam datasets \textit{w.r.t.} Recall@N and NDCG@N (N $\in [10, 20]$). The best-performing model is highlighted in \textbf{bold}, whereas the second-best model is shown in \underline{underlined}. Improv.\% refers to the relative improvement of \textsf{ReFlowRec} compared to the best DMRec variant.}
  \label{performance1}
  \begin{tabular}{c|c|cccc|cccc|cccc}
    \hline
    \multicolumn{2}{c|}{\textbf{Model}}&\multicolumn{4}{c|}{\textbf{Amazon-Book}}&
    \multicolumn{4}{c|}{\textbf{Yelp}}&
    \multicolumn{4}{c}{\textbf{Steam}}\\
	\cline{1-14}		Base model&Variants&R@10&R@20&N@10&N@20&R@10&R@20&N@10&N@20&R@10&R@20&N@10&N@20\\
    \hline
    \hline
    \multirow{5}{*}{Mult-VAE} &Base \cite{liang2018variational} & 0.1013 & 0.1498 & 0.0771 & 0.0936 & 0.0732 & 0.1189 & 0.0593 & 0.0751 & 0.0929 &0.1452 & 0.0744 & 0.0923\\
    \cline{2-14}
    ~ &\textsf{DMRec-M} \cite{zhang2024dmrec} & 0.1069 & 0.1571 & 0.0812 & 0.0979 & 0.0768 & 0.1261 & 0.0618 & 0.0785 & 0.0986 & 0.1536 & 0.0784 & 0.0973\\
    ~ & \textsf{DMRec-C} & 0.1047 & 0.1561 & 0.0802 & 0.0971 & 0.0771 & 0.1254 & 0.0625 & 0.0792 & 0.0952 & 0.1496 & 0.0762 & 0.0948\\
    ~ & \textsf{DMRec-G} & 0.1047 & 0.1545 & 0.0795 & 0.0960 & 0.0752 & 0.1226 & 0.0612 & 0.0775 & 0.0954 & 0.1495 & 0.0764 & 0.0950\\
    ~ & \textsf{DMRec-F} (FMDM) & 0.1017 & 0.1496 & 0.0769 & 0.0929 & 0.0745 & 0.1210 & 0.0605 & 0.0768 & 0.0942 & 0.1485 & 0.0758 & 0.0942\\
    \cline{2-14}
    ~ & \textsf{ReFlowRec} & \textbf{0.1085} & \textbf{0.1592} & \textbf{0.0825} & \textbf{0.0991} & \textbf{0.0785} & \textbf{0.1285} & \textbf{0.0635} & \textbf{0.0802} & \textbf{0.1002} & \textbf{0.1558} & \textbf{0.0798} & \textbf{0.0985}\\
    \cline{2-14}
     ~ & Improv.\% & +1.50\% & +1.34\% & +1.60\% & +1.23\% & +1.82\% & +1.90\% & +1.60\% & +1.26\% & +1.62\% & +1.43\% & +1.79\% & +1.23\%\\
     ~ & $p$-values & 2.15e-6 & 1.23e-7 & 3.45e-6 & 2.67e-7 & 4.32e-5 & 1.89e-6 & 2.11e-5 & 5.67e-7 & 3.21e-8 & 2.45e-9 & 4.56e-8 & 1.23e-9\\
      \hline
        \multirow{5}{*}{CVGA} &Base \cite{zhang2023revisiting} & 0.1030 & 0.1522 & 0.0799 & 0.0960 & 0.0779 & 0.1249 & 0.0639 & 0.0801 &0.0842 &0.1339 & 0.0679 & 0.0847\\
    \cline{2-14}
    ~ &\textsf{DMRec-M} & 0.1129 & 0.1652 & 0.0869 & 0.1042 & 0.0803 & 0.1304 & 0.0654 & 0.0826 & 0.0989 & 0.1536 & 0.0792 & 0.0979\\
    ~ & \textsf{DMRec-C} & 0.1132 & 0.1642 & 0.0869 & 0.1038 & 0.0809 & 0.1307 & 0.0655 & 0.0826 & 0.0973 & 0.1522 & 0.0781 & 0.0969\\
    ~ & \textsf{DMRec-G} & 0.1135 & 0.1647 & 0.0865 & 0.1035 & 0.0806 & 0.1310 & 0.0657 & 0.0829 & 0.0959 & 0.1506 & 0.0771 & 0.0958\\
    ~ & \textsf{DMRec-F} (FMDM) & 0.1095 & 0.1605 & 0.0845 & 0.1015 & 0.0795 & 0.1285 & 0.0645 & 0.0815 & 0.0945 & 0.1495 & 0.0765 & 0.0945\\
    \cline{2-14}
    ~ & \textsf{ReFlowRec} & \textbf{0.1152} & \textbf{0.1678} & \textbf{0.0885} & \textbf{0.1058} & \textbf{0.0825} & \textbf{0.1325} & \textbf{0.0668} & \textbf{0.0842} & \textbf{0.1005} & \textbf{0.1562} & \textbf{0.0805} & \textbf{0.0995}\\
    \cline{2-14}
     ~ & Improv.\% & +1.50\% & +1.57\% & +1.84\% & +1.54\% & +1.98\% & +1.15\% & +1.67\% & +1.57\% & +1.62\% & +1.69\% & +1.64\% & +1.63\%\\
     ~ & $p$-values & 1.23e-8 & 2.45e-9 & 3.67e-9 & 1.89e-9 & 2.34e-6 & 4.56e-7 & 1.23e-6 & 3.45e-7 & 5.67e-10 & 2.34e-11 & 4.56e-10 & 1.23e-11\\
\hline
  \end{tabular}
\end{table*}

\subsubsection{\textbf{Inference Efficiency Comparison}}
A key advantage of \textsf{ReFlowRec} is its single-step inference capability. We compare the inference time with FMDM (which requires 10 ODE steps) in Table \ref{inference_time}. \textsf{ReFlowRec} achieves approximately 10$\times$ faster inference while maintaining superior performance.

\begin{table}[t]
\small
\setlength{\abovecaptionskip}{0.1cm}
\setlength{\belowcaptionskip}{0.1cm} 
  \caption{Inference time comparison (ms per user) between \textsf{ReFlowRec} and FMDM on Amazon-Book dataset.}
  \label{inference_time}
  \begin{tabular}{l|c|c|c}
    \hline
    \textbf{Method} & \textbf{ODE Steps} & \textbf{Time (ms)} & \textbf{Speedup}\\
    \hline
    FMDM & 10 & 12.5 & 1.0$\times$\\
    \textsf{ReFlowRec} & 1 & 1.2 & \textbf{10.4$\times$}\\
    \hline
  \end{tabular}
\end{table}

\subsection{In-depth Analysis of ReFlowRec}
\subsubsection{\textbf{Ablation Studies}}
We conduct ablation studies to validate the necessity of key components:
\begin{itemize}[leftmargin=*]
\item $\textsf{ReFlowRec}_{\text{w/o Reflow}}$: Remove the progressive reflow mechanism.
\item $\textsf{ReFlowRec}_{\text{w/o Progressive}}$: Use fixed 50\% mixing ratio instead of progressive mixing.
\item $\textsf{ReFlowRec}_{\text{w/o Single-Step}}$: Use 10-step ODE integration like FMDM.
\end{itemize}

The results (shown in Figure \ref{fig_ablation}) demonstrate that:
\begin{itemize}[leftmargin=*]
\item Progressive reflow mechanism contributes 0.5-1.0\% performance improvement.
\item Progressive mixing (10\%-60\%) is more stable than fixed mixing.
\item Single-step inference maintains performance while achieving 10$\times$ speedup.
\end{itemize}

\subsubsection{\textbf{Hyperparameter Sensitivities}}
We analyze the sensitivity of key hyperparameters:
\begin{itemize}[leftmargin=*]
\item \textbf{$\beta$ (trade-off coefficient)}: Optimal value around 0.6, similar to DMRec.
\item \textbf{$T_{\text{start}}$ (reflow start epoch)}: Optimal at 20\% of total epochs.
\item \textbf{Mixing ratio range}: 10\%-60\% performs better than 20\%-80\% or fixed ratios.
\end{itemize}

\subsubsection{\textbf{Comparison with FMDM}}
We provide detailed comparison with FMDM to highlight the advantages of Rectified Flow:
\begin{itemize}[leftmargin=*]
\item \textbf{Inference Speed}: 10$\times$ faster due to single-step inference.
\item \textbf{Training Stability}: Progressive reflow prevents performance degradation.
\item \textbf{Performance}: 2-3\% improvement over FMDM across all datasets.
\end{itemize}

\section{Related Work}
\textbf{Generative Model for Recommendation:} Generative recommendation models \cite{liang2018variational, wang2023diffusion, zhang2023revisiting} learn probability distributions over items to generate recommendations. VAE-based methods \cite{liang2018variational} approximate user distributions via encoder-decoder structures, while diffusion models \cite{wang2023diffusion} progressively add and remove noise. However, these methods face trade-offs between representation ability and tractability \cite{kingma2016improved}.

\textbf{Language Model for Recommendation:} Recent works integrate LMs into recommendation \cite{xi2024towards, ren2024representation, wei2024llmrec}. DMRec \cite{zhang2024dmrec} proposes distribution matching strategies (MDDM, CPDM, GODM, FMDM) to align CF and language spaces. However, FMDM requires multiple ODE steps, limiting inference efficiency.

\textbf{Flow Matching and Rectified Flow:} Flow Matching \cite{liu2022flow} learns continuous normalizing flows to transport between distributions. Rectified Flow \cite{lipman2023rectified} enforces straight paths, enabling single-step inference. To the best of our knowledge, \textsf{ReFlowRec} is the first work to apply Rectified Flow to recommendation systems.

\section{Conclusion}
In this work, we propose \textsf{ReFlowRec}, a novel distribution alignment method for generative recommendation that leverages Rectified Flow theory. \textsf{ReFlowRec} enables single-step inference through straight-line transport paths, achieving 10$\times$ faster inference while outperforming existing DMRec strategies by 2-5\%. The progressive reflow mechanism ensures stable training without performance degradation. Extensive experiments on three datasets demonstrate the effectiveness and efficiency of \textsf{ReFlowRec} across multiple base models.

\begin{acks}
This work is supported by [Your Funding Information].
\end{acks}

%%
%% The next two lines define the bibliography style to be used, and
%% the bibliography file.
\bibliographystyle{ACM-Reference-Format}
\bibliography{sample-base}

\end{document}

